\worksheettitle
  {Домашняя работа: ответы}
  {\modulemark{M01-H} Версия учителя}

\begin{teacherbox}[Назначение]
  Домашняя работа закрепляет только результат M01: ученик определяет разряд
  цифры и называет её значение. Разложение числа на разрядные слагаемые здесь
  не проверяется, поскольку относится к M02.
\end{teacherbox}

\section{Ответы}

\begin{enumerate}[label=\textbf{\arabic*.}]
  \item \(6\); \(3\); \(1\).

  \item
    \begin{itemize}
      \item в числе \(5\,832\) цифра \(8\) стоит в разряде сотен и имеет
        значение \(800\);
      \item в числе \(46\,075\) цифра \(6\) стоит в разряде единиц тысяч и
        имеет значение \(6\,000\);
      \item в числе \(90\,204\) цифра \(2\) стоит в разряде сотен и имеет
        значение \(200\).
    \end{itemize}

  \item \(70\,000\), \(700\), \(7\). Значения различаются, потому что цифры
    занимают разные разряды.

  \item Цифра \(6\) стоит в разряде единиц тысяч, поэтому её значение равно
    \(6\,000\).
\end{enumerate}

\begin{resultbox}[Как оценивать]
  \begin{itemize}
    \item \textbf{Результат достигнут:} задания 2--4 выполнены верно, в
      объяснении названа зависимость значения цифры от разряда.
    \item \textbf{Нужна M01-R:} ученик называет саму цифру вместо её значения,
      определяет разряды слева направо или не различает значения одинаковых
      цифр.
  \end{itemize}
\end{resultbox}
